% GENERATED FILE. Edit canonical lessons, then run npm run generate:books.
% canonical-source-hash: fnv1a64:6e526637d24ca8ef
% canonical-lessons: RU-C534-budit, RU-C534-vysypatsya, RU-C534-prichyosyvatsya, RU-C534-krasitsya, RU-C534-khalat, RU-R534-review-533-534, RU-R534-review-final-534

\chapter{Mornings at Home}
\label{ch:ru-mornings-at-home}
\hlchaptermodality{534}

\begin{chapteropening}
\textbf{By the end of this chapter:} \emph{I can describe getting ready in the morning.}

Complete the last lesson of chapter 534: I can describe getting ready in the morning.
\end{chapteropening}

\section[\texorpdfstring{\ru{будить} (budít)}{будить (budít)}]{\ru{будить} (budít) — to wake (someone)}

\label{lesson:RU-C534-budit}



\begin{quote}
Before the new one: say the Russian for to burn out, then the Russian for insulating tape.
\end{quote}

\subsection*{You'll want to know: \ru{будить}}
\textbf{\ru{будить}} (\emph{budít}) — \textquotedblleft{}to wake (someone)\textquotedblright{}.

\begin{quote}
\ru{Каждое} \ru{утро} \ru{мама} \ru{будит} \ru{меня} \ru{в} \ru{семь}.
\end{quote}

\begin{quote}
\emph{Kázhdoye útro máma búdit menyá v sem.}
\end{quote}

\begin{quote}
Every morning Mum wakes me at seven.
\end{quote}

Imperfective; present \ru{бужу}, \ru{будишь} (\ru{д} becomes \ru{ж}). Perfective \ru{разбудить}.

\begin{cousinweb}
\emph{\ru{будить}} is inherited from Proto-Slavic \emph{buditi} 'to awaken', which EtymDB traces to PIE \emph{b\textsuperscript{h}ewd\textsuperscript{h}- 'to awaken'.}
\end{cousinweb}

\begin{grammarlens}[title={What you've built}]
One more word for telling, explaining and describing at B1.
\end{grammarlens}

\subsection*{Your turn}
\begin{itemize}
\raggedright
  \item \emph{Say it:} \emph{budít}
  \item \emph{Say it:} \emph{budít}, once more
  \item \emph{Say it:} a sentence of your own with \emph{budít}
  \item \emph{Recall:} say the Russian for to burn out, then the Russian for insulating tape, then say \emph{budít} again
\end{itemize}

\begin{tcolorbox}[breakable,colback=teal!4,colframe=teal!35!black,title={Before you move on}]
What does \ru{будить} mean? (\textquotedblleft{}To wake (someone)\textquotedblright{}.) Say it once more.
\end{tcolorbox}
\section[\texorpdfstring{\ru{высыпаться} (vysypátsya)}{высыпаться (vysypátsya)}]{\ru{высыпаться} (vysypátsya) — to get enough sleep}

\label{lesson:RU-C534-vysypatsya}



\begin{quote}
Before the new one: say the Russian for insulating tape, then the Russian for to wake.
\end{quote}

\subsection*{You'll want to know: \ru{высыпаться}}
\textbf{\ru{высыпаться}} (\emph{vysypátsya}) — \textquotedblleft{}to get enough sleep\textquotedblright{}.

\begin{quote}
\ru{В} \ru{выходные} \ru{я} \ru{наконец} \ru{высыпаюсь}.
\end{quote}

\begin{quote}
\emph{V vykhodnýye ya nakonéts vysypáyus.}
\end{quote}

\begin{quote}
At the weekend I finally get enough sleep.
\end{quote}

Imperfective; perfective \ru{выспаться} (\ru{высплюсь}), with the stress on \emph{\ru{вы}-}.

\begin{cousinweb}
\emph{\ru{высыпаться}} is built on \emph{\ru{вы}-} and the root of \emph{\ru{спать}} 'to sleep', with \emph{-\ru{ся}}; EtymDB has no entry.
\end{cousinweb}

\begin{grammarlens}[title={What you've built}]
One more word for telling, explaining and describing at B1.
\end{grammarlens}

\subsection*{Your turn}
\begin{itemize}
\raggedright
  \item \emph{Say it:} \emph{vysypátsya}
  \item \emph{Say it:} \emph{vysypátsya}, once more
  \item \emph{Say it:} a sentence of your own with \emph{vysypátsya}
  \item \emph{Recall:} say the Russian for insulating tape, then the Russian for to wake, then say \emph{vysypátsya} again
\end{itemize}

\begin{tcolorbox}[breakable,colback=teal!4,colframe=teal!35!black,title={Before you move on}]
What does \ru{высыпаться} mean? (\textquotedblleft{}To get enough sleep\textquotedblright{}.) Say it once more.
\end{tcolorbox}
\section[\texorpdfstring{\ru{причёсываться} (prichyósyvatsya)}{причёсываться (prichyósyvatsya)}]{\ru{причёсываться} (prichyósyvatsya) — to do one's hair}

\label{lesson:RU-C534-prichyosyvatsya}



\begin{quote}
Before the new one: say the Russian for to wake, then the Russian for to get enough sleep.
\end{quote}

\subsection*{You'll want to know: \ru{причёсываться}}
\textbf{\ru{причёсываться}} (\emph{prichyósyvatsya}) — \textquotedblleft{}to do one's hair\textquotedblright{}.

\begin{quote}
\ru{Сестра} \ru{долго} \ru{причёсывается} \ru{перед} \ru{зеркалом}.
\end{quote}

\begin{quote}
\emph{Sestrá dólgo prichyósyvayetsya péred zérkalom.}
\end{quote}

\begin{quote}
My sister spends a long time doing her hair in front of the mirror.
\end{quote}

Imperfective; perfective \ru{причесаться} (\ru{причешусь}).

\begin{cousinweb}
\emph{\ru{причёсываться}} is built on \emph{\ru{при}-} and \emph{\ru{чесать}} 'to comb', which is inherited from Proto-Slavic \emph{česati} (EtymDB).
\end{cousinweb}

\begin{grammarlens}[title={What you've built}]
One more word for telling, explaining and describing at B1.
\end{grammarlens}

\subsection*{Your turn}
\begin{itemize}
\raggedright
  \item \emph{Say it:} \emph{prichyósyvatsya}
  \item \emph{Say it:} \emph{prichyósyvatsya}, once more
  \item \emph{Say it:} a sentence of your own with \emph{prichyósyvatsya}
  \item \emph{Recall:} say the Russian for to wake, then the Russian for to get enough sleep, then say \emph{prichyósyvatsya} again
\end{itemize}

\begin{tcolorbox}[breakable,colback=teal!4,colframe=teal!35!black,title={Before you move on}]
What does \ru{причёсываться} mean? (\textquotedblleft{}To do one's hair\textquotedblright{}.) Say it once more.
\end{tcolorbox}
\section[\texorpdfstring{\ru{краситься} (krásitsya)}{краситься (krásitsya)}]{\ru{краситься} (krásitsya) — to put on make-up}

\label{lesson:RU-C534-krasitsya}



\begin{quote}
Before the new one: say the Russian for to get enough sleep, then the Russian for to do one's hair.
\end{quote}

\subsection*{You'll want to know: \ru{краситься}}
\textbf{\ru{краситься}} (\emph{krásitsya}) — \textquotedblleft{}to put on make-up\textquotedblright{}.

\begin{quote}
\ru{Она} \ru{не} \ru{красится}, \ru{когда} \ru{сидит} \ru{дома}.
\end{quote}

\begin{quote}
\emph{Oná ne krásitsya, kogdá sidít dóma.}
\end{quote}

\begin{quote}
She doesn't wear make-up when she's at home.
\end{quote}

Imperfective; perfective \ru{накраситься}. Also to dye one's hair. The reflexive of \emph{\ru{красить}} — to paint.

\begin{cousinweb}
\emph{\ru{краситься}} is the reflexive of \emph{\ru{красить}} 'to paint', built on \emph{\ru{краска}} 'paint', from the root of \emph{\ru{краса}} 'beauty'; EtymDB has no entry for the reflexive.
\end{cousinweb}

\begin{grammarlens}[title={What you've built}]
One more word for telling, explaining and describing at B1.
\end{grammarlens}

\subsection*{Your turn}
\begin{itemize}
\raggedright
  \item \emph{Say it:} \emph{krásitsya}
  \item \emph{Say it:} \emph{krásitsya}, once more
  \item \emph{Say it:} a sentence of your own with \emph{krásitsya}
  \item \emph{Recall:} say the Russian for to get enough sleep, then the Russian for to do one's hair, then say \emph{krásitsya} again
\end{itemize}

\begin{tcolorbox}[breakable,colback=teal!4,colframe=teal!35!black,title={Before you move on}]
What does \ru{краситься} mean? (\textquotedblleft{}To put on make-up\textquotedblright{}.) Say it once more.
\end{tcolorbox}
\section[\texorpdfstring{\ru{халат} (khalát)}{халат (khalát)}]{\ru{халат} (khalát) — a dressing gown; a white coat}

\label{lesson:RU-C534-khalat}



\begin{quote}
Before the new one: say the Russian for to do one's hair, then the Russian for to put on make-up.
\end{quote}

\subsection*{You'll want to know: \ru{халат}}
\textbf{\ru{халат}} (\emph{khalát}) — \textquotedblleft{}a dressing gown; a white coat\textquotedblright{}.

\begin{quote}
\ru{Утром} \ru{папа} \ru{пьёт} \ru{кофе} \ru{в} \ru{халате}.
\end{quote}

\begin{quote}
\emph{Útrom pápa pyot kófe v khaláte.}
\end{quote}

\begin{quote}
In the morning Dad drinks his coffee in his dressing gown.
\end{quote}

Masculine. Also a doctor's or a cook's white coat.

\begin{cousinweb}
EtymDB lists \emph{\ru{халат}}, glossed 'dressing gown; robe', without a source; it is usually traced through a Turkic language to Arabic \emph{xilʿa} 'a robe of honour'.
\end{cousinweb}

\begin{grammarlens}[title={What you've built}]
One more word for telling, explaining and describing at B1.
\end{grammarlens}

\subsection*{Your turn}
\begin{itemize}
\raggedright
  \item \emph{Say it:} \emph{khalát}
  \item \emph{Say it:} \emph{khalát}, once more
  \item \emph{Say it:} a sentence of your own with \emph{khalát}
  \item \emph{Recall:} say the Russian for to do one's hair, then the Russian for to put on make-up, then say \emph{khalát} again
\end{itemize}

\begin{tcolorbox}[breakable,colback=teal!4,colframe=teal!35!black,title={Before you move on}]
What does \ru{халат} mean? (\textquotedblleft{}A dressing gown; a white coat\textquotedblright{}.) Say it once more.
\end{tcolorbox}
\section[(dialogue)]{Leisure, the arts and home — a review}

\label{lesson:RU-R534-review-533-534}



\begin{quote}
Say the Russian for a screw and pliers.
\end{quote}

\subsection*{Your turn}
Say these aloud:

\begin{itemize}
\raggedright
  \item a screw, pliers, to unscrew, to burn out, insulating tape
  \item to wake, to get enough sleep, to do one's hair, to put on make-up, a dressing gown
\end{itemize}

\begin{tcolorbox}[breakable,colback=teal!4,colframe=teal!35!black,title={Before you move on}]
A screw? (\textbf{\ru{шуруп}}.) Pliers? (\textbf{\ru{плоскогубцы}}.) To unscrew? (\textbf{\ru{откручивать}}.) To burn out? (\textbf{\ru{перегорать}}.) Insulating tape? (\textbf{\ru{изолента}}.) To wake? (\textbf{\ru{будить}}.) To get enough sleep? (\textbf{\ru{высыпаться}}.) To do one's hair? (\textbf{\ru{причёсываться}}.) To put on make-up? (\textbf{\ru{краситься}}.) A dressing gown? (\textbf{\ru{халат}}.)
\end{tcolorbox}
\section[(dialogue)]{Mornings at Home — once more}

\label{lesson:RU-R534-review-final-534}



\begin{quote}
Say the Russian for to wake and to get enough sleep.
\end{quote}

\subsection*{Your turn}
\begin{itemize}
\raggedright
  \item \emph{Say it:} to wake, to get enough sleep, to do one's hair, to put on make-up, a dressing gown
\end{itemize}

\begin{tcolorbox}[breakable,colback=teal!4,colframe=teal!35!black,title={Before you move on}]
To wake? (\textbf{\ru{будить}}.) To get enough sleep? (\textbf{\ru{высыпаться}}.) To do one's hair? (\textbf{\ru{причёсываться}}.) To put on make-up? (\textbf{\ru{краситься}}.) A dressing gown? (\textbf{\ru{халат}}.)
\end{tcolorbox}
